% 05-signatures.tex: cryptographic evidence.
\section{Authenticating truth: signatures, custody, two witnesses and one migration}
\label{sec:signatures}

\motif{An entry in the register is not proof of who made it.}

\analogy{Every entry in the register carries the seal of the official who made it, pressed with a die nobody else holds. Two assayers, trained in different schools, examine each new seal, and a seal that only one of them recognises is not accepted.}

\subsection{What a signature is for}

A digital signature is a seal that only the holder of a private key can make and that anyone with
the matching public key can check. Polaris signs the SHA3-256 digest of a token value under the
issuing agency's ML-DSA key and stores the signature beside the public key that made it. The claim
a signature carries is narrow and permanent: \emph{this key signed these bytes.} It says nothing
about whether the token is still authorized, which is the next section's problem, and nothing about
who holds the key, which is custody's problem. \sImpl

The algorithm is ML-DSA-65 (FIPS 204, security category 3) by default and ML-DSA-87 (category 5)
beside it; ML-DSA-44 is refused everywhere as below the floor. No signed body hardcodes its
algorithm: the custodied key decides it before signing, because \code{algorithm} is itself a signed
field, and the registry advertises the accepted set. SLH-DSA-128s and SLH-DSA-256s, the hash-based
alternatives that do not rest on lattice assumptions, are registered in the algorithm table
precisely because they rest on hash functions alone, so that a rotation away from lattices is a row
update; no SLH-DSA signer is wired, and the posture ledger carries the gap. \sImpl{} for ML-DSA;
\sFut{} for a hash-based signer.

\subsection{Custody}

The private key never sits in application code. It sits behind a custody interface with three
drivers: a file (development), PKCS\#11 (a hardware security module or a software token; the key is
generated inside the token and never leaves it), and a cloud key-management service. Continuous
integration exercises the PKCS\#11 driver against a software token that implements ML-DSA, not a
hardware module, and rotates a key inside it: a token signed under the old in-token key still
verifies after the switch while the new key signs. The KMS driver's envelope cryptography is real
and its service is a stand-in. One production profile makes the module the sole signer: the
application refuses to boot if a file key is present or real signing is unavailable. Which driver,
which module, and who holds the key are decisions the readiness ledger assigns to a deploying
organisation. \sImpl{} for the drivers and the drill; \sExt{} for hardware custody, since no hardware
module has been used.

The placeholder deserves a paragraph, because it is where an honest system is most easily
misread. Without the real-signing flag and a
lattice library present, the signing module produces a deterministic SHA3-256 binding labelled
\code{DETERMINISTIC-PLACEHOLDER-SHA3-256}: thirty-two bytes that verify against no key, where a real
ML-DSA-65 signature is 3,309. That is the default in development and in most of continuous
integration. It is guarded: a production boot fails closed unless real signing is
actually available, the placeholder is a named development profile, and it warns loudly when used
unnamed, and it is still the reason the repository no longer calls itself post-quantum without
saying so in the same breath. \sImpl

\subsection{Two witnesses at issuance, one at use, sampled, and one outside corpus}

The two-witness principle states that any cryptographic verdict Polaris ships must be checkable by
a second, independent implementation, written separately, differently represented, sharing no code,
and agreeing on the verdict or abstaining explicitly. A verifier that agrees with itself has
established nothing. At issuance the signature just produced is verified by two independent FIPS 204
implementations, a lattice library and an OpenSSL-backed one, and is refused persistence unless both
accept it; a missing witness at issuance is a refusal, not a downgrade. \sImpl

Verification at use runs one witness, because a stored signature is already known to verify under
both, and one witness detects tampering at about ten times the rate: on the reference machine
roughly 7,950 verifications per second per core against roughly 740 two-witness. The fast path
earns its speed by being continuously checked: a mandatory random fraction of successful checks is
replayed through the second witness, a disagreement increments a counter and pages at the highest
severity, and every response names which witness set ran. \sImpl

One thing about the witnesses is not self-referential. On every push, the Project Wycheproof
ML-DSA-65 verification vectors, a third-party corpus pinned to a commit in which accepting a
signature the corpus calls invalid is a failure, run under both witnesses: 58 of 58 agree. That is
evidence that Polaris's verification of the primitive agrees with an outside authority on 58 cases.
It is not a wallet, not an interoperability result, and it says nothing about the credential
protocol; the repository's scoreboard records it in exactly those words. \sWit{} for the primitive.

The same principle governs the zero-knowledge verdict (Section~\ref{sec:privacy}): a Python
re-implementation with plain integers modulo a prime, sharing no code with the Rust prover,
re-derives the membership statement and must agree, and it abstains, in the open, on the one axis it
cannot model, the integrity of the proof bytes. What the principle catches is implementation
divergence. What it does not catch is a shared misreading of the specification, and the repository
says so beside every two-witness claim. The witnesses are also mutation-tested, so
the claim that two implementations must agree is a claim about the tests as well as the code.

\subsection{Canonical signing: the same bytes on both sides}

Every signed artifact in Polaris is a JSON statement with sorted keys and no whitespace, digested
with SHA3-256, then signed. The application builds the statement; the detached verifier
(Section~\ref{sec:verifiers}) rebuilds it from the artifact to check the signature. The two
constructions must be byte-identical or a genuine artifact fails to verify, and a test, the
canonical-equivalence oracle, pins every signed type on both sides. The one bug the oracle is there
to catch was found in review: an artifact whose \code{algorithm} field was set after signing carried
a signed field the signature did not cover. \sImpl

\subsection{Keys have a lifecycle}

The key register is an append-only table of events, registered, retired, or compromised from an
instant that may predate the discovery, from which a signed trust list is published. A verifier
holding a trust list decides a key's status at any instant, so a credential under a compromised
issuer key is rejected independently of the issuer's own manifest, which a thief holding the key
could forge, and a document signed and timestamped before the compromise stays valid while one
signed after does not. The two-instance drill records a compromise on one authority and watches the
other's decision flip over the wire, with no cooperation from the compromised party. \sImpl

\subsection{Agility under an audited migration path, and what a break would cost}

ML-DSA-65's security does not rest on a prediction about quantum hardware; it rests on Module-LWE and Module-SIS, the lattice problems the scheme is built on, being hard
against the best known classical and quantum algorithms, which could be wrong for mathematical
reasons. The two real risks are a cryptanalytic advance against the lattice assumptions and an
implementation flaw, and the design already answers both, the hash-based hedge in the registry and
the two witnesses at issuance. What the repository can demonstrate, and so what it claims, is \emph{algorithm agility under an audited migration path}: C7 makes the
algorithm a row, \entity{AgencyAlgorithmAuth} authorizes it per authority, and a population can be
re-signed under a new parameter set on a path that runs on every push. \sImpl

That path is measured rather than planned. \code{polaris migrate-population} re-signs the whole
active population under a new parameter set, built out of the same constraints the per-token
procedure relies on: the unique constraint is the serialization point, so two runners racing one
credential produce one row, and sixty-four runners divide the population without coordinating. The
number that matters is not throughput but how many holders hold a credential that verifies under
nothing, and the drill samples it at every batch boundary rather than at the ends, because a gap
that opens and closes between two endpoints is invisible to a before-and-after check. It is zero
throughout. The old signature keeps verifying until its deprecation date, and that interval
\emph{is} the migration: deprecating the superseded algorithm is refused while any active credential
is unmigrated, because a holder would otherwise discover the gap at a border while the console
reported progress. Under real ML-DSA-87 with one custodied key the drill re-signs about 345
credentials per second on one runner, 91 percent of it signing, so a population of 350 million is
about twelve days on one runner and four and a half hours on sixty-four; buying database capacity for
it buys almost nothing. \sImpl

What a break in Module-LWE would cost is written down in the readiness ledger, and this paper repeats
it. It would not require a schema change, a change on the verification path, or a new trust model.
It would cost every credential already issued under the broken algorithm, and one thing no agility
buys back: during a cutover a credential carries a classical signature beside the post-quantum one,
and the classical half is protected by nothing this design provides, so a credential harvested today
is readable on the day its classical signature falls, whatever the system migrates to afterwards.
That asymmetry is why signing post-quantum now is worth doing, and it is the limit of what doing it
achieves.

The repository's own lab then asked whether the agility is on the side that needs it, and found that
it is issuer-side only. The algorithm's \code{deprecation\_date} reaches no signed artifact: the
registry publishes a bare list of algorithm names, the trust list carries per-key status only, and
the detached verifier's accepted set is a dictionary in code. So a verifier cannot learn that an
algorithm has been deprecated, retiring one is a software release to every relying party rather than
a migration, and key revocation does not substitute, because if the assumption breaks an attacker
forges under any key the trust list calls active. Rollback and the mixed window during a partial
migration are unmeasured. No change is proposed; the finding is recorded as a limitation for a
deploying organisation. \sExt

\subsection{What remains classical}

The token signature, the digests, the Merkle commitments, the zero-knowledge proof, the
client-to-edge key exchange for modern clients and the hop from the
application to the connection pooler are post-quantum, each of the last two read off a real
handshake rather than inferred from a base image. The hop from the pooler to the database, the certificate
signatures, and the operator's WebAuthn credentials (hardware-key sign-in) remain classical, each gated on a third party (a
PostgreSQL release, the public certificate ecosystem, authenticator hardware) and each stated in the
posture ledger with its threat and its exposure. None of the lattice library, the proof system or the
circuit has had an external cryptographic audit. \sExt

\begin{layercard}{Layer card: cryptographic evidence}
\qa{What it is}{ML-DSA signatures over SHA3-256 digests, stored with their public keys; a custody interface with file, PKCS\#11 and KMS drivers; two independent verifying implementations and one outside corpus; a canonical statement format; an append-only key register and a signed trust list; a population migration that runs.}
\qa{Why it exists}{A claim that leaves the database needs to carry its own proof, a proof only one program can check is a promise, and an algorithm that cannot be replaced is a deadline.}
\qa{What it trusts}{FIPS 204 and 202 as standardised; the two libraries not sharing a bug; the custody backend the operator chose.}
\qa{What it refuses}{A verdict from a single implementation at issuance; an algorithm named in code rather than by the key; a signature stored without its key; a placeholder passing as a signature; a deprecation while any credential would be left verifying under nothing.}
\qa{What it can see}{The bytes it signs, which are digests: the signing path never sees a person.}
\qa{What it cannot see}{Whether the signed claim is still authorized (Section~\ref{sec:authorization}); whether a verifier in the field knows an algorithm has been retired.}
\qa{If it fails}{A forged signature fails both witnesses; a compromised key is retired from an instant by the register and every verifier holding the trust list rejects what it signed afterwards; a broken algorithm costs every credential under it and is survived by re-signing, at a measured rate.}
\qa{Checked by}{\code{check\_pqc\_signing\_wired}, \code{check\_pqc\_second\_witness}, \code{check\_verify\_witness\_sampling}, \code{check\_real\_pqc\_default\_boot}, \code{check\_trust\_lifecycle}, \code{check\_algorithm\_agility}, \code{check\_quantum\_event\_readiness}, \code{check\_zk\_witnesses\_are\_mutation\_tested}, \code{check\_post\_quantum\_claims\_are\_agility}, the canonical-equivalence oracle, the attack suite.}
\qa{Now or later}{Now. The hardware module and the external audit are later, and belong to a deploying organisation; the signed algorithm status a verifier could read is not built.}
\qa{Inside and outside}{Inside: the schema, which stores the signature rows. Outside: the independent verifiers that check them without the issuer.}
\end{layercard}

\nextlayer{A signature the issuer's own server checks proves nothing to a stranger, because the
stranger has to trust the server. The next layer takes the verification out of the issuer's hands
entirely, and, for the first time, into the hands of software the issuer's author did not write.}
